% ============================================================
% Capítulo 11 — Evaluación: conceptos + código
% ============================================================
\chapter{Evaluación: extracción de features y métricas}
\label{cap:evaluacion}

Este capítulo explica los conceptos de inferencia y muestra el
código de evaluación. Para la explicación detallada de las
métricas (Top-k, mAP, AP), ver el \cref{cap:metricas}.

\section{¿Qué es la inferencia (evaluación)?}

Una vez entrenado el modelo, queremos medir su rendimiento. El
proceso tiene tres pasos:

\begin{enumerate}
  \item \textbf{Extraer embeddings} de todas las imágenes y
        textos del split test.
  \item \textbf{Computar la matriz de similitud coseno} entre
        imágenes y textos.
  \item \textbf{Calcular métricas} (Top-k, mAP) usando los pids.
\end{enumerate}

\noindent\textbf{Diferencia con entrenamiento}:
\begin{itemize}
  \item En entrenamiento se hace forward + backward.
  \item En inferencia sólo forward (sin gradientes).
\end{itemize}

\section{¿Por qué \texttt{@torch.no\_grad()}?}

\texttt{@torch.no\_grad()} es un decorator que desactiva el
cálculo de gradientes. Beneficios:

\begin{itemize}
  \item \textbf{Menor memoria}: no se almacenan gradientes.
  \item \textbf{Más rápido}: ahorra cómputo.
  \item \textbf{Previene errores}: garantiza que no se
        modifiquen parámetros accidentalmente.
\end{itemize}

\section{¿Qué es \texttt{model.eval()}?}

\texttt{model.eval()} cambia el modo del modelo:

\begin{itemize}
  \item Desactiva \texttt{Dropout} (no se eliminan activaciones).
  \item Usa estadísticas \emph{acumuladas} en \texttt{BatchNorm}
        en lugar de las del batch.
  \item Afecta a capas como \texttt{Dropout}, \texttt{BatchNorm}.
\end{itemize}

\noindent Para inferir SIEMPRE se debe llamar
\texttt{model.eval()} antes.

\section{Galería vs query: convención de este proyecto}

En este proyecto:

\begin{itemize}
  \item \textbf{Query} = texto (descripción de la persona).
  \item \textbf{Gallery} = imágenes.
\end{itemize}

\noindent Para cada query se buscan imágenes similares en la
gallery. La métrica \textbf{Top-k} indica cuántas queries
tienen la imagen correcta en el top-k.

\section{Código: extracción de features de imágenes}

\begin{minted}[language=Python, numbers, firstnumber=42, linerange=42-62]{python}
@torch.no_grad()
def extract_image_features(model, img_loader, device):
    """Extrae embeddings visuales para todas las imágenes del loader."""

    # Modo evaluación (desactiva dropout, etc.)
    model.eval()

    all_features = []
    all_pids = []
    all_paths = []

    for batch in tqdm(img_loader, desc='Extracting image features'):
        if isinstance(batch, (list, tuple)):
            imgs, pids, paths = batch[0], batch[1], batch[2]
        else:
            imgs, pids, paths = batch

        # GPU
        imgs = imgs.to(device)

        # Forward solo en la rama visual
        # image_embedding retorna (B, 1024, 1)
        feats = model.image_embedding(imgs)

        # Aplanar y mover a CPU
        feats = feats.view(feats.size(0), -1).cpu()

        all_features.append(feats)
        all_pids.extend(pids.tolist() if torch.is_tensor(pids) else list(pids))
        all_paths.extend(paths if isinstance(paths, list) else [paths])

    # Concatenar todos los batches
    return torch.cat(all_features, 0), all_pids, all_paths
\end{minted}

\section{Código: extracción de features de textos}

\begin{minted}[language=Python, numbers, firstnumber=65, linerange=65-85]{python}
@torch.no_grad()
def extract_text_features(model, txt_loader, device):
    """Extrae embeddings textuales."""

    model.eval()

    all_features = []
    all_pids = []

    for batch in tqdm(txt_loader, desc='Extracting text features'):
        if isinstance(batch, (list, tuple)):
            pids, token_ids, token_lengths = batch[0], batch[1], batch[2]
        else:
            pids, token_ids, token_lengths = batch

        token_ids = token_ids.to(device)
        token_lengths = token_lengths.to(device)

        # Forward solo en la rama textual
        feats = model.text_embedding(token_ids, token_lengths)
        feats = feats.view(feats.size(0), -1).cpu()

        all_features.append(feats)
        all_pids.extend(pids.tolist() if torch.is_tensor(pids) else list(pids))

    return torch.cat(all_features, 0), all_pids
\end{minted}

\section{Código: similitud coseno}

\begin{minted}[language=Python, numbers, firstnumber=9, linerange=9-17]{python}
def calculate_similarity(image_feature_local, text_feature_local):
    """image: (N_img, 1024), text: (N_txt, 1024)
       returns: (N_img, N_txt) numpy array de similitudes coseno"""

    # Aplanar por si tienen más dims
    image_feature_local = image_feature_local.view(image_feature_local.size(0), -1)
    text_feature_local = text_feature_local.view(text_feature_local.size(0), -1)

    # Normalizar L2 (cada vector a norma 1)
    image_feature_local = image_feature_local / \
        image_feature_local.norm(dim=1, keepdim=True)
    text_feature_local = text_feature_local / \
        text_feature_local.norm(dim=1, keepdim=True)

    # Producto matricial: S[i,j] = cos(img_i, txt_j)
    similarity = torch.mm(image_feature_local, text_feature_local.t())

    # Devolver como numpy (requerido por calculate_ap)
    return similarity.cpu().numpy()
\end{minted}

\section{Código: pipeline completo de evaluación}

\begin{minted}[language=Python, numbers, firstnumber=89, linerange=89-153]{python}
def main():
    args = parse_args()

    config = sys_configuration(
        dataset_name="CUHK-PEDES",
        dataset_source=args.dataset_source
    )
    config['batch_size'] = args.batch_size
    config['device'] = args.device
    config['model_testing_data_split'] = args.split  # 'test' o 'val'

    # Construir dataloaders (devuelve img_loader y txt_loader)
    _, img_loader, txt_loader, _ = build_cuhkpedes_dataloader(config)

    # Construir modelo y cargar checkpoint
    model = TextReIDNet(config).to(config.device)
    ckpt = torch.load(args.checkpoint, map_location=config.device)
    if 'model_state_dict' in ckpt:
        model.load_state_dict(ckpt['model_state_dict'])
    else:
        model.load_state_dict(ckpt)  # compatibilidad

    # PASO 1: extraer features de imágenes
    img_feats, img_pids, img_paths = extract_image_features(
        model, img_loader, config.device
    )

    # PASO 2: extraer features de textos
    txt_feats, txt_pids = extract_text_features(
        model, txt_loader, config.device
    )

    # PASO 3: matriz de similitud coseno
    similarity = calculate_similarity(img_feats, txt_feats)

    # PASO 4: calcular CMC y mAP
    img_pids_t = torch.tensor(img_pids)
    txt_pids_t = torch.tensor(txt_pids)
    cmc, mAP = evaluate(similarity, txt_pids_t, img_pids_t)

    # PASO 5: imprimir resultados
    top1  = cmc[0] * 100
    top5  = cmc[4] * 100 if len(cmc) > 4 else 0
    top10 = cmc[9] * 100 if len(cmc) > 9 else 0

    print(f"  Top-1:  {top1:.2f}%")
    print(f"  Top-5:  {top5:.2f}%")
    print(f"  Top-10: {top10:.2f}%")
    print(f"  mAP:    {mAP*100:.2f}%")

    # Guardar en log
    logger = setup_logger(name='eval_logger',
                           log_file_path=config.test_log_path,
                           write_mode='append')
    logger.info(f"Top-1: {top1:.2f}% | Top-5: {top5:.2f}% | "
                f"Top-10: {top10:.2f}% | mAP: {mAP*100:.2f}%")
\end{minted}

\section{Uso desde la terminal}

\begin{minted}[language=bash, basicstyle=\ttfamily\small, frame=single]{python}
# Evaluar en el split test
.venv/bin/python evaluate.py \
    --checkpoint data/checkpoints/TextReIDNet_latest.pth.tar \
    --dataset_source huggingface \
    --split test

# Cambiar batch size
.venv/bin/python evaluate.py \
    --checkpoint data/checkpoints/TextReIDNet_latest.pth.tar \
    --batch_size 64

# Forzar CPU (lento)
.venv/bin/python evaluate.py \
    --checkpoint data/checkpoints/TextReIDNet_latest.pth.tar \
    --device cpu
\end{minted}

\section{Resultados típicos}

\begin{table}[H]
\centering
\caption{Resultados esperados según el README.}
\label{tab:resultados-tipicos-eval}
\begin{tabular}{lcccc}
\toprule
\textbf{Setup} & \textbf{Top-1} & \textbf{Top-5} & \textbf{Top-10} & \textbf{mAP} \\
\midrule
Sin entrenar (random) & \textasciitilde 0.02\% & 0.05\% & 0.10\% & 0.10\% \\
HF (128×128) 50 batches & 0.02\% & 0.08\% & 0.17\% & 0.10\% \\
HF entrenado completo & 45--52\% & \textasciitilde 70\% & \textasciitilde 80\% & \textasciitilde 40\% \\
Original paper & 54.02\% & 75.51\% & 85.28\% & \textasciitilde 50\% \\
\bottomrule
\end{tabular}
\end{table}

\section{Limitaciones}

\begin{enumerate}
  \item \textbf{Sin gallery augmentation}: cada imagen se evalúa
        una sola vez.
  \item \textbf{Sin query expansion}: no se re-rankea usando
        los top-k.
  \item \textbf{Sin análisis por longitud de caption}.
  \item \textbf{Sin visualización} de aciertos/fallos.
\end{enumerate}
